\textbf{COLUNI 2025 -- Questão 35}\par\medskip Durante o período da Era Vargas (1930-1945), o Brasil passou por transformações políticas, econômicas e sociais. Sobre essas transformações e o processo que levou Vargas ao poder, considere as afirmativas a seguir:
\begin{enumerate}
\item[I.] Vargas chegou ao poder com o apoio das elites paulistas e dos militares de alta patente, cujo discurso coincidia com o da Aliança Liberal na defesa de um Brasil moderno e industrial. Isso fortaleceu o governo durante todo o período de sua vigência.
\item[II.] Durante a década de 1920, a Semana de Arte Moderna e o Tenentismo colocaram em xeque a formação cultural e política do Brasil. Cada um, à sua maneira, defendeu uma modernidade pautada na valorização do elemento nacional.
\item[III.] Durante a década de 1930, Vargas combateu os comunistas e os fascistas no Brasil. Em 1935, na chamada “intentona comunista” Vargas perseguiu e prendeu tenentes e membros da Aliança Nacional Libertadora (ANL). Em 1938, a Ação Integralista Brasileira (AIB) foi fechada e teve seus membros perseguidos pelo Estado Novo.
\item[IV.] Nas décadas de 1930 e 1940, o governo Vargas investiu pesado na indústria nacional, canalizando parte das riquezas do café para novos setores produtivos, como os bens de consumo duráveis. Na indústria de base, o Estado fundou a CSN (Companhia Siderúrgica Nacional) e a CVRD (Companhia Vale do Rio Doce) com parcerias e empréstimos estrangeiros, aumentando a dívida externa.
\item[V.] Nos centros urbanos, a classe operária se aproximou de Vargas. Durante o Estado Novo (1937-1945), com os direitos trabalhistas e o controle da mídia nacional, o governo usou o DIP (Departamento de Imprensa e Propaganda) para criar uma imagem positiva do ditador que ficou conhecido como “pai dos pobres”.
\end{enumerate}
Está CORRETO o que se afirma apenas em:
\begin{enumerate}
\item II, III e V.
\item II, IV e V.
\item I, II e III.
\item III, IV e V.
\end{enumerate}
